\LTR

\textbf{Answer briefly. Unless a calculation is explicitly requested, each sub-question should be answered in at most 1--2 lines.}

\vspace{1.5em}
\noindent \textbf{Q1. Human Visual System (HVS) Sensitivity}
\begin{itemize}
    \item (a) The HVS is most sensitive to which spatial frequencies?
    \begin{itemize}
        \item[$\circ$] A) Very low only
        \item[$\circ$] B) Very high only
        \item[$\circ$] C) Mid-frequencies (approx. 3-10 cycles/degree)
        \item[$\circ$] D) All frequencies equally
    \end{itemize}
    \item (b) Is the eye generally more sensitive to \textbf{horizontal/vertical} gratings or to \textbf{diagonal} gratings?
\end{itemize}

\vspace{1.5em}
\noindent \textbf{Q2. Image Fidelity Metrics}\\
A reconstructed image has slightly altered pixel values but preserves all main edges and structural details. Which metric is more likely to align with human visual judgment: \textbf{MSE}, \textbf{PSNR}, or \textbf{SSIM}? Provide a one-line justification.

\vspace{1.5em}
\noindent \textbf{Q3. 1-D Aliasing Intuition}\\
A sinusoid with a frequency of \textbf{8 Hz} is sampled at a rate of \textbf{8.5 Hz}. Briefly explain:
\begin{itemize}
    \item (a) Why aliasing occurs in this scenario.
    \item (b) Why the resulting sampled points might suggest a much lower-frequency waveform than the original.
\end{itemize}

\vspace{1.5em}
\noindent \textbf{Q4. Quantization Concepts}\\
\hspace*{1.5em}\textbf{Give one concise reason for each part.}
\begin{itemize}
    \item (a) Why does a \textbf{Lloyd-Max quantizer} outperform a uniform quantizer when the input values are not equally likely (e.g., Gaussian distribution)?
    \item (b) Why is a \textbf{Compandor} (Compressor-Expander) often preferred in practical hardware over the direct implementation of a fully non-uniform quantizer?
\end{itemize}

\RTL